% Part: computability
% Chapter: recursive-functions
% Section: other-recursions

\documentclass[../../../include/open-logic-section]{subfiles}

\begin{document}

\olfileid{cmp}{rec}{ore}
\olsection{عوديات أخرى}

وبالاقتران والجمع في متتاليات نرد إلى العودية البدائية صورًا أخرى،
أبعد عن صورتها الأولى وأوسع نفعًا. فمن ذلك تعريف دالتين معًا،
على الوجه الآتي:
\begin{align*}
{\OLId{latin-ordinary}{h}}_{\OLMathNumeral{0}}(\vec {\OLId{latin-ordinary}{x}}, {\OLMathNumeral{0}}) & = {\OLId{latin-ordinary}{f}}_{\OLMathNumeral{0}}(\vec {\OLId{latin-ordinary}{x}}) \\
{\OLId{latin-ordinary}{h}}_{\OLMathNumeral{1}}(\vec {\OLId{latin-ordinary}{x}}, {\OLMathNumeral{0}}) & = {\OLId{latin-ordinary}{f}}_{\OLMathNumeral{1}}(\vec {\OLId{latin-ordinary}{x}}) \\
{\OLId{latin-ordinary}{h}}_{\OLMathNumeral{0}}(\vec {\OLId{latin-ordinary}{x}}, {\OLId{latin-ordinary}{y}}+{\OLMathNumeral{1}}) & = {\OLId{latin-ordinary}{g}}_{\OLMathNumeral{0}}(\vec {\OLId{latin-ordinary}{x}}, {\OLId{latin-ordinary}{y}}, {\OLId{latin-ordinary}{h}}_{\OLMathNumeral{0}}(\vec {\OLId{latin-ordinary}{x}}, {\OLId{latin-ordinary}{y}}), {\OLId{latin-ordinary}{h}}_{\OLMathNumeral{1}}(\vec {\OLId{latin-ordinary}{x}}, {\OLId{latin-ordinary}{y}})) \\
{\OLId{latin-ordinary}{h}}_{\OLMathNumeral{1}}(\vec {\OLId{latin-ordinary}{x}}, {\OLId{latin-ordinary}{y}}+{\OLMathNumeral{1}}) & = {\OLId{latin-ordinary}{g}}_{\OLMathNumeral{1}}(\vec {\OLId{latin-ordinary}{x}}, {\OLId{latin-ordinary}{y}}, {\OLId{latin-ordinary}{h}}_{\OLMathNumeral{0}}(\vec {\OLId{latin-ordinary}{x}}, {\OLId{latin-ordinary}{y}}), {\OLId{latin-ordinary}{h}}_{\OLMathNumeral{1}}(\vec {\OLId{latin-ordinary}{x}}, {\OLId{latin-ordinary}{y}}))
\end{align*}
وهذا ما يسمى \emph{العودية المتزامنة}. ومن وجوه التعريف النافعة أيضًا
أن نعيّن ${\OLId{latin-ordinary}{h}}(\vec {\OLId{latin-ordinary}{x}},{\OLId{latin-ordinary}{y}}+{\OLMathNumeral{1}})$ بدلالة \emph{جميع} القيم
${\OLId{latin-ordinary}{h}}(\vec {\OLId{latin-ordinary}{x}},{\OLMathNumeral{0}})$، …،~${\OLId{latin-ordinary}{h}}(\vec {\OLId{latin-ordinary}{x}},{\OLId{latin-ordinary}{y}})$، كما في التعريف الآتي:
\begin{align*}
{\OLId{latin-ordinary}{h}}(\vec {\OLId{latin-ordinary}{x}}, {\OLMathNumeral{0}}) & = {\OLId{latin-ordinary}{f}}(\vec {\OLId{latin-ordinary}{x}}) \\
{\OLId{latin-ordinary}{h}}(\vec {\OLId{latin-ordinary}{x}}, {\OLId{latin-ordinary}{y}}+{\OLMathNumeral{1}}) & = {\OLId{latin-ordinary}{g}}(\vec {\OLId{latin-ordinary}{x}}, {\OLId{latin-ordinary}{y}}, \tuple{{\OLId{latin-ordinary}{h}}(\vec {\OLId{latin-ordinary}{x}}, {\OLMathNumeral{0}}), \dots, {\OLId{latin-ordinary}{h}}(\vec {\OLId{latin-ordinary}{x}}, {\OLId{latin-ordinary}{y}})}).
\end{align*}
ويجمع المخطط الآتي هذا المعنى في عبارة أقصر:
\[
{\OLId{latin-ordinary}{h}}(\vec {\OLId{latin-ordinary}{x}}, {\OLId{latin-ordinary}{y}}) = {\OLId{latin-ordinary}{g}}(\vec {\OLId{latin-ordinary}{x}}, {\OLId{latin-ordinary}{y}}, \tuple{{\OLId{latin-ordinary}{h}}(\vec {\OLId{latin-ordinary}{x}}, {\OLMathNumeral{0}}), \dots, {\OLId{latin-ordinary}{h}}(\vec {\OLId{latin-ordinary}{x}}, {\OLId{latin-ordinary}{y}}-{\OLMathNumeral{1}})})
\]
على أن تكون الحجة الأخيرة لـ${\OLId{latin-ordinary}{g}}$ المتتالية الخالية عند
${\OLId{latin-ordinary}{y}}={\OLMathNumeral{0}}$. والمعنى في الصورتين واحد: للدالة ${\OLId{latin-ordinary}{h}}$ عند حساب
«خطوة اللاحق» أن تستعمل سائر القيم التي حُسبت قبلها، مجموعة في متتالية.
وهذا هو \emph{عودية مسار القيم}. ومن تطبيقاته
جواز تعريف من القبيل الآتي:
\begin{align*}
{\OLId{latin-ordinary}{h}}(\vec {\OLId{latin-ordinary}{x}}, {\OLId{latin-ordinary}{y}}) & = \begin{cases}
  {\OLId{latin-ordinary}{g}}(\vec {\OLId{latin-ordinary}{x}}, {\OLId{latin-ordinary}{y}}, {\OLId{latin-ordinary}{h}}(\vec {\OLId{latin-ordinary}{x}}, {\OLId{latin-ordinary}{k}}(\vec {\OLId{latin-ordinary}{x}}, {\OLId{latin-ordinary}{y}}))) & \text{إذا كان ${\OLId{latin-ordinary}{k}}(\vec {\OLId{latin-ordinary}{x}}, {\OLId{latin-ordinary}{y}}) < {\OLId{latin-ordinary}{y}}$} \\
  {\OLId{latin-ordinary}{f}}(\vec {\OLId{latin-ordinary}{x}}) & \text{في غير ذلك}
\end{cases}
\end{align*}
فتُحسب قيمة ${\OLId{latin-ordinary}{h}}$ عند ${\OLId{latin-ordinary}{y}}$ من قيمة ${\OLId{latin-ordinary}{h}}$ عند
\emph{أي} حجة سابقة، تتولى~${\OLId{latin-ordinary}{k}}$ تعيينها.

\begin{prob}
  ضع لدالة الباقي ${\OLId{latin-ordinary}{r}}({\OLId{latin-ordinary}{x}},{\OLId{latin-ordinary}{y}})$ تعريفًا بعودية مسار القيم. (إذا كان ${\OLId{latin-ordinary}{x}}$ و${\OLId{latin-ordinary}{y}}$
  عددين طبيعيين وكان ${\OLId{latin-ordinary}{y}}>{\OLMathNumeral{0}}$، فإن ${\OLId{latin-ordinary}{r}}({\OLId{latin-ordinary}{x}},{\OLId{latin-ordinary}{y}})$ هو العدد الأصغر من~${\OLId{latin-ordinary}{y}}$ الذي
  يحقق ${\OLId{latin-ordinary}{x}}={\OLId{latin-ordinary}{z}}\times {\OLId{latin-ordinary}{y}}+{\OLId{latin-ordinary}{r}}({\OLId{latin-ordinary}{x}},{\OLId{latin-ordinary}{y}})$ لبعض~${\OLId{latin-ordinary}{z}}$. وللتعيين، لنقل إنه إذا كان
  ${\OLId{latin-ordinary}{y}}={\OLMathNumeral{0}}$ فإن ${\OLId{latin-ordinary}{r}}({\OLId{latin-ordinary}{x}},{\OLMathNumeral{0}})={\OLMathNumeral{0}}$.)
\end{prob}

فتأمل كيف تُستخرج هذه الدوال بالعودية البدائية
المعتادة. وبقيت صورة أوسع تصرفًا، نجيز فيها
تبديل \emph{المعلمات}، أي القيم الجانبية، أثناء العودية:
\begin{align*}
{\OLId{latin-ordinary}{h}}(\vec {\OLId{latin-ordinary}{x}}, {\OLMathNumeral{0}}) & = {\OLId{latin-ordinary}{f}}(\vec {\OLId{latin-ordinary}{x}}) \\
{\OLId{latin-ordinary}{h}}(\vec {\OLId{latin-ordinary}{x}}, {\OLId{latin-ordinary}{y}}+{\OLMathNumeral{1}}) & = {\OLId{latin-ordinary}{g}}(\vec {\OLId{latin-ordinary}{x}}, {\OLId{latin-ordinary}{y}}, {\OLId{latin-ordinary}{h}}({\OLId{latin-ordinary}{k}}(\vec {\OLId{latin-ordinary}{x}}), {\OLId{latin-ordinary}{y}}))
\end{align*}
وهذه أيضًا تحاكيها العودية البدائية المعتادة. وردّها إليها يحتاج إلى دقة؛
ومن سبيل الاهتداء إليه أن تبسط الحساب بيدك خطوة خطوة.

\end{document}
